\documentclass[11pt]{article}
\usepackage{amsmath,amssymb,amsthm}
\usepackage[margin=1in]{geometry}
\newtheorem{theorem}{Theorem}
\newtheorem{proposition}[theorem]{Proposition}
\newtheorem{lemma}[theorem]{Lemma}
\newtheorem{corollary}[theorem]{Corollary}
\theoremstyle{remark}
\newtheorem{remark}[theorem]{Remark}
\begin{document}
\title{Problem 1, Cell 6: $N$ lines in $\mathbb R^d$ (partial)}
\date{26 September 2026}
\maketitle

\noindent\fbox{\parbox{\dimexpr\linewidth-2\fboxsep-2\fboxrule}{\textbf{Official statement (Cell 6, open
question).} ``Fejes T\'oth's conjecture from the Setting, for every $N$ and every $d$. For $N$ lines in
$\mathbb R^d$, write $N=qd+s$ with $0\le s<d$, and let $M(N,d)=s\binom{q+1}2+(d-s)\binom q2$. Prove or
disprove: every $N$ lines in $\mathbb R^d$ satisfy $S\le\bigl(\binom N2-M(N,d)\bigr)\cdot\pi/2$. This is
the value attained by splitting the lines as evenly as possible among $d$ mutually orthogonal
directions. Settling any infinite family not already covered above counts as partial progress.''}}
\medskip

\paragraph{Status: Partial.} \emph{Proved:}
(i)~the conjecture for $N=6$ lines in $\mathbb R^3$ (Theorem~\ref{thm:63}), by averaging Cell~4/5;
(ii)~the exact list of what subset averaging can give from the cases proved so far. For $d\ge3$ and
$N\ge d+3$ it is the single pair $(6,3)$, so averaging yields \emph{no} new infinite family
(Proposition~\ref{prop:avg});
(iii)~a structural reduction of the case $N=d+3$ for all $d$. At a saturated minimiser, every
component of the non-orthogonality graph is settled, except components of rank $\ge4$, nullity
$\ge3$, maximum degree $3$ and cyclomatic number $\ge2$ (Theorem~\ref{thm:red}).
\emph{Not proved:} the conjecture in general, and the case $N=d+3$ for any $d\ge4$. The first open
instance is $7$ lines in $\mathbb R^4$. We do \emph{not} claim a new infinite family.

\section{Setting and known cases}
Lines are given by unit vectors $x_1,\dots,x_N\in\mathbb R^d$, $g_{ij}=\langle x_i,x_j\rangle$,
\[
A(x)=\sum_{i<j}\arcsin|g_{ij}|,\qquad S=\binom N2\frac\pi2-A .
\]
So the conjecture reads $A\ge M(N,d)\,\pi/2$. For earlier work on the problem see \cite{FVZ16,BM19}. Writing $n_1,\dots,n_d$ for the sizes of an even split,
$M(N,d)=\sum_i\binom{n_i}2$. Hence $M(N,d)\ge N-d$ (since $\binom n2\ge n-1$ for integers
$n\ge0$), and $M(N,d)=\max(0,N-d)$ for $N\le 2d$. Also $A$ only increases when lines are added, since all
terms are $\ge0$. The following cases are proved in our earlier cells and used here.
\begin{itemize}
\item[(C1)] $d=2$, all $N$: $A\ge M(N,2)\pi/2$.
\item[(C3)] $N=d+1$, all $d\ge1$: $A\ge\pi/2$.
\item[(C5)] $N=d+2$, all $d\ge1$: $A\ge\pi$. This contains C4 (five lines in $\mathbb R^3$, six in
  $\mathbb R^4$).
\end{itemize}
From C5 we also use the orthogonality-saturation proposition, the path lemma, the block-structure
lemma and the cycle lemma (C5, Prop.~1 and Lemmas~2--4). These give a minimiser of $A$ in which each
vertex of the non-orthogonality graph $J$ ($ij\in J$ iff $g_{ij}\ne0$) is either isolated or has
orthogonal partners spanning $\ge d-1$ dimensions. For $N=d+k$: every vertex of $J$ has degree
$\le k$. Distinct components $C$ of $J$ span orthogonal subspaces of dimensions $r_C$ with
$\sum_Cr_C\le d$, and $A=\sum_CA_C$. A vertex of $C$ of degree $\delta\ge1$ forces
$r_C\le m_C-\delta$, where $m_C=|C|$. Write $\nu_C=m_C-r_C$ (the \emph{nullity} of $C$).

\section{Subset averaging}

\begin{lemma}[Averaging]\label{lem:avg}
If every $n$ lines in $\mathbb R^d$ satisfy $A\ge B\,\pi/2$, then every $N\ge n$ lines in
$\mathbb R^d$ satisfy
\[
A\ \ge\ \frac{N(N-1)}{n(n-1)}\,B\,\frac\pi2 .
\]
\end{lemma}
\begin{proof}
Sum $A$ over the $\binom Nn$ subsets of size $n$. Each pair $\{i,j\}$ lies in $\binom{N-2}{n-2}$
of them, so $\binom{N-2}{n-2}A\ge\binom NnB\frac\pi2$, and
$\binom Nn/\binom{N-2}{n-2}=\frac{N(N-1)}{n(n-1)}$.
\end{proof}

\begin{theorem}\label{thm:63}
Any $6$ lines in $\mathbb R^3$ satisfy $A\ge3\pi/2=M(6,3)\,\pi/2$, i.e.\ $S\le12\cdot\frac\pi2=6\pi$.
Equality holds for three orthogonal directions, each taken twice.
\end{theorem}
\begin{proof}
By C4/C5, every $5$ lines in $\mathbb R^3$ have $A\ge\pi$. There are $\binom65=6$ five-subsets, and each
pair lies in $\binom43=4$ of them. So $4A\ge6\pi$, and $A\ge\frac32\pi$. Here $M(6,3)$ has $q=2$, $s=0$,
so $M=3\binom22=3$. Also $\binom62-3=12$. For three orthogonal doubled directions, the three
coincident pairs give $A=3\cdot\frac\pi2$.
\end{proof}

\begin{proposition}[Averaging gives no infinite family]\label{prop:avg}
Let $d\ge3$ and $N\ge d+3$. Averaging (Lemma~\ref{lem:avg}) from any of the cases proved in
\textup{C3}, \textup{C5} or Theorem~\ref{thm:63} gives $A\ge M(N,d)\pi/2$ only for $(N,d)=(6,3)$.
\end{proposition}
\begin{proof}
For $d\ge3$ the available sources are $n\le d$ ($B=0$, useless), $n=d+1$ ($B=1$), $n=d+2$ ($B=2$),
and $(6,3)$ with $B=3$. C1 concerns $\mathbb R^2$ and does not apply to lines spanning more. The
averaged bound is $\frac{N(N-1)}{n(n-1)}B$. Its coefficient $B/(n(n-1))$ is
$\frac1{d(d+1)}$, $\frac2{(d+1)(d+2)}$ and (for $d=3$) $\frac3{30}=\frac2{20}$. For $d\ge2$ the second
is the largest ($2d\ge d+2$). So we must decide when
\[
\beta(N,d):=\frac{2N(N-1)}{(d+1)(d+2)}\ \ge\ M(N,d).
\]
\emph{Range $d+3\le N\le2d$.} Here $M=N-d$. Put $N=d+k$, $3\le k\le d$. The condition is
$f(k):=2(d+k)(d+k-1)-k(d+1)(d+2)\ge0$, and $f$ is convex in $k$. We have
$f(3)=(d+2)(3-d)$ and $f(d)=d\bigl(8d-4-(d+1)(d+2)\bigr)=-d(d-2)(d-3)$. For $d\ge4$ both are
negative, so $f<0$ on $[3,d]$. For $d=3$ only $k=3$ occurs, and $f(3)=0$: this is $(6,3)$, with
equality.

\emph{Range $N\ge2d+1$.} By convexity $M(N,d)\ge d\binom{N/d}2=\frac{N(N-d)}{2d}$. It suffices that
$4d(N-1)<(d+1)(d+2)(N-d)$, i.e.\ $(d^2-d+2)N>d(d+1)(d+2)-4d$. The left side increases with $N$. At
$N=2d+1$ the difference of the two sides is $d^3-4d^2+5d+2=d(d-2)^2+d+2>0$. So $\beta<M$ for all
$N\ge2d+1$.
\end{proof}

\begin{remark}
Averaging can also use sources with a different $n$ at once. This does not help, because the averaged
bound from the case $(N',d)$ just obtained is at most the bound from the original source (the factors
$\frac{N(N-1)}{n(n-1)}$ multiply). The script \texttt{averaging\_search.py} checks
Proposition~\ref{prop:avg} in exact arithmetic for $3\le d\le60$, $d+3\le N\le400$. It prints the
single pair $(6,3)$.
\end{remark}

\section{The case $N=d+3$: reduction to hard components}

\begin{lemma}[PSD forcing]\label{lem:force}
Let $C$ be a connected graph on vertex set $[m]$, and let $x_1,\dots,x_m$ be unit vectors with
$\langle x_u,x_v\rangle\ne0$ iff $uv\in E(C)$ ($u\ne v$). Let $\nu=m-\dim\operatorname{span}(x_v)$.
If $C$ is a tree, then $\nu\le1$. If $C$ has exactly one cycle, then $\nu\le2$.
\end{lemma}
\begin{proof}
$\nu$ is the dimension of $Z=\{c\in\mathbb R^m:\sum_vc_vx_v=0\}$. Call a set $B$ \emph{forcing} if every
$c\in Z$ vanishing on $B$ vanishes everywhere; then $c\mapsto c|_B$ is injective on $Z$, so
$\nu\le|B|$. \emph{Forcing rule.} Let $c\in Z$ vanish on $B$, and let $K$ be a component of $C-B$.
Vectors in different components of $C-B$ are orthogonal. So the vectors
$\sum_{v\in K}c_vx_v$ (one for each $K$) are pairwise orthogonal and add up to $0$, hence each is $0$. If
$w\in B$ has exactly one neighbour $u$ in $K$, then
$0=\langle x_w,\sum_{v\in K}c_vx_v\rangle=c_ug_{wu}$, so $c_u=0$ and $u$ may be added to $B$.

\emph{Tree:} $B=\{\text{root}\}$. A vertex of $B$ has exactly one neighbour (its child) in each
subtree hanging below it, so its children are forced, and by induction every vertex is.
\emph{One cycle} $v_1\cdots v_\ell$ ($\ell\ge3$) with trees attached: $B=\{v_1,v_2\}$. Removing $B$, the
vertex $v_\ell$ lies in a component in which $v_1$ has no other neighbour. So $v_\ell$ is forced, then
$v_{\ell-1}$ by $v_\ell$, and so on around the cycle. The attached trees are forced as in the tree
case. Hence $\nu\le|B|$ in both cases.
\end{proof}

\begin{theorem}[Reduction for $N=d+3$]\label{thm:red}
Let $x$ be a saturated minimiser of $A$ for $d+3$ lines in $\mathbb R^d$ (C5, Prop.~1). Then
$A\ge3\pi/2$, unless $J$ has a component $C$ with $r_C\ge4$, $\nu_C\ge3$, a vertex of degree $3$
(and maximum degree $3$), and cyclomatic number $|E(C)|-m_C+1\ge2$.
\end{theorem}
\begin{proof}
$\sum_C\nu_C=N-\sum_Cr_C\ge3$. We show $A_C\ge\min(\nu_C,3)\frac\pi2$ for every component that is not of
the excluded type. Since $\min(a,3)+\min(b,3)\ge\min(a+b,3)$, this gives $A\ge\frac32\pi$ whenever no
excluded component exists.

\emph{$\nu_C\le2$.} The $m_C=r_C+\nu_C$ lines lie in $V_C\cong\mathbb R^{r_C}$. By C3 ($\nu_C=1$) and C5
($\nu_C=2$), $A_C\ge\nu_C\frac\pi2$.

\emph{$r_C\le3$.} If $r_C\le2$: C1 in $\mathbb R^{r_C}$ (or coincident lines when $r_C=1$) gives
$A_C\ge M(m_C,r_C)\frac\pi2\ge(m_C-r_C)\frac\pi2=\nu_C\frac\pi2$. If $r_C=3$ and $\nu_C\ge3$: any $6$ of the
$m_C\ge6$ lines satisfy Theorem~\ref{thm:63}, and $A_C\ge\frac32\pi$ by monotonicity.

\emph{The rest}: $r_C\ge4$ and $\nu_C\ge3$. Degrees are $\le3$ (block lemma, $k=3$). If $C$ had maximum
degree $\le2$, it would be an edge, a path or a cycle. The C5 analysis then gives $\nu_C\le2$ (a path on
$\ge3$ vertices is impossible; a triangle or cycle has $\nu=2$). So $C$ has a vertex of degree $3$. By
Lemma~\ref{lem:force}, $\nu_C\ge3$ excludes trees and unicyclic graphs, so the cyclomatic number is
$\ge2$. This is exactly the excluded type.
\end{proof}

\begin{corollary}
The conjecture holds for $N=d+3$ when $d\le3$, and for every configuration of $d+3$ lines whose
saturated minimiser has no excluded component. Since $\sum r_C\le d$ and an excluded component needs
$r_C\ge4$ and $m_C\ge7$, the first open case is $7$ lines in $\mathbb R^4$ with $J$ connected.
\end{corollary}

For $d\le3$ this is not new: $d=2$ is C1 and $d=3$ is Theorem~\ref{thm:63}. The corollary settles no
new infinite family.

\section{Why the general case stays open}
For the excluded components there is no analogue of the cycle lemma. It would have to be an inequality
$A_C\ge\frac32\pi$ for vector families whose non-orthogonality graph has cyclomatic number $\ge2$ and
maximum degree $3$. The script \texttt{scan7.py} looks at the first open case, $7$ lines in
$\mathbb R^4$. It lists the $41$ connected graphs on $7$ vertices with maximum degree $\le3$ and
cyclomatic number $\ge2$. For each, it minimises $A$ with the orthogonality pattern imposed (penalty
method, $4$ random starts). In \emph{every} pattern the infimum found is $1.5\pi=M(7,4)\frac\pi2$. It is
approached only as some prescribed non-zero inner products tend to $0$, i.e.\ at the boundary of the
pattern, where the block structure changes. So no interior local minimum with connected $J$ appeared
below $\frac32\pi$. This is consistent with the conjecture, but it is evidence only. A proof would need
a potential (like $\Phi$ in C5) for graphs with two independent cycles.

\begin{remark}[What we could not reproduce]
Our working notes mentioned a rank-$4$, $7$-line configuration with $A\approx1.702\pi$ as a witness that
the hard components are genuinely different. The file describing it (\texttt{p1\_c6\_codex.md}) was not
available when this cell was written, and we did not reconstruct it. We therefore make no claim about it.
\end{remark}

\section*{Scripts}
\begin{itemize}
\item \texttt{averaging\_search.py}: Proposition~\ref{prop:avg}, exact arithmetic, $\approx3$\,s.
\item \texttt{scan7.py}: the pattern scan for $7$ lines in $\mathbb R^4$ (numerical, not a proof),
  a few minutes.
\item \texttt{search7.py}: unconstrained local minimisation for $7$ lines in $\mathbb R^4$ ($150$ random
  starts). Every run ends at $1.5\pi$, with $J$ split into components of sizes $(1,2,2,2)$ or
  $(2,2,3)$.
\end{itemize}

\begin{thebibliography}{9}
\bibitem{FVZ16} F.~Fodor, V.~V\'igh, T.~Zarn\'ocz, On the angle sum of lines, \emph{Arch.\ Math.} 106
(2016) 91--100.
\bibitem{BM19} D.~Bilyk, R.~Matzke, On the Fejes T\'oth problem about the sum of angles between lines,
\emph{Proc.\ Amer.\ Math.\ Soc.} 147 (2019), arXiv:1801.07837.
\end{thebibliography}
\end{document}
